Take the wholly prepaid invoice in a synthetic book containing only cash 30
and a contract liability of 30. No performance has occurred. Advance A1 has
30 unused; there is no invoice yet. The new invoice I1 records gross billings
of 30 and applies all of A1. Cumulative $G$ and $K$ both rise from zero to 30,
so $T=G+A-K$ remains 30. Cash, revenue, receivables and contract balances do
not change. But A1 is now fully used and I1 is paid. Omitting this transition
would let a later invoice claim the same credit again.

The authored advance-allocation profile therefore pins a closed journal
variant and a separate economic delta. In addition to their common context,
basis and authorization references, their distinguishing fields are:
\begin{lstlisting}
JournalProposal {
  effect: "state_only", lines: []
}
EconomicDelta {
  account_changes: zero for every chart account,
  before_contract: advance A1 unused 30; no invoice,
  after_contract: A1 used 30; invoice I1 gross 30,
                  applied 30, unpaid 0; allocation A1->I1
}
\end{lstlisting}
This display abbreviates the required common references; the pinned schema
bundle supplies the closed field lists and complete portable records. The
independent calculation checks both levels. It derives zero account changes
from the contract equations and verifies that the advance, invoice and
allocation agree with the new cumulative values. Empty lines with a proposed
receivable increase of 1 fail that reconciliation. Zero account changes with
the allocation omitted also fail: the required economic effect is missing.

Authorization uses the identified unused advance, a fresh invoice identity,
the zero unpaid amount and the admitted contract bounds. A caller cannot
replace those findings with a true permission flag. The profile supports a
synthetic book only; its historical source support and authored facts provide
no authorization for a real company's posting. An assurance that the empty
journal balances would establish none of these conditions.

For a valid first event E1, the atomic transaction publishes the state-only
proposal, independent delta, new contract state, receipt and event-index entry,
and advances the branch generation once. The receipt is evidence that the
allocation was committed, even though every account balance stayed the same.
The stable identity remains the entity, book, branch and business event.
Its fingerprint covers the whole portable request, including the original
expected head, generation, request identifier and execution settings.

Retrying the identical request after the commit returns the original receipt.
The lookup must precede comparison with today's head; otherwise a successful
retry would wrongly fail because its own earlier commit advanced the branch.
Changing I1 to I2 under E1 instead conflicts, as does changing the request
identifier while keeping the amounts. A genuinely new event E2 submitted
against the old opening fails freshness. Rebuilding E2 against the current
opening cannot rescue the reused credit: A1 has no unused amount, so
authorization fails before lookup. These are different reasons for preserving
the current book unchanged.

On replay, the response identifies the original committed event but reports
no additional commit by this invocation. A lost acknowledgement remains
unknown until lookup under the original identity resolves it; inventing a new
event identity could double the allocation. Exports and recovery must also
retain the economic transition. A format limited to monetary journal lines
must preserve the state-only event separately or report that a lossless export
is unsupported. Dropping it because revenue and cash did not change would
make the recovered advance available twice. The complete documentary exchanges
exercise these distinctions; transaction durability and concurrent execution
remain implementation checks.
